\subsection{Fixed finite edge covers do not yield all-prime exactness}

\begin{proposition}[Proved: a persistent obstruction for edge covers]
\label{C4:voltage-cover-exclusion}
Let $I$ be a nonempty finite ordered set and let $\sigma,\tau$
be permutations of $I$. For $A\in\Cdelta$, define
$F_{I,\sigma,\tau}(A)$ on the lexicographically ordered set
$A\times I$ by replacing each increasing edge $a<b<c$ by
\[
 \bigl\{(a,i),(b,\sigma(i)),(c,\tau(i))\bigr\},\qquad i\in I.
 \tag{COV.1}
\]
On arrows put $F(f)(v,i)=(f(v),i)$. This is a strong
endofunctor of $\Cdelta$.

There is one fixed three-arrow family in $\Cdelta$ whose
image under every such functor, and under every finite
composition of such functors, has no affine equalizer over
any nonzero unital ring. Consequently a filtered colimit of
copies of $\Cdelta$ with these transition functors remains
nonexact over every nonzero unital ring, although it is total.
\end{proposition}
\begin{proof}
For $U\subseteq F(A)$ write
$m_v=|U\cap(\{v\}\times I)|$ and
$A_t=\{v:m_v\geq t\}$. Above an old edge $a<b<c$, at
most $\min(m_a,m_b,m_c)$ lifted edges lie in $U$, since
all three fiber maps in (COV.1) are injective. Therefore
\[
 \delta(U)\geq \sum_{t=1}^{|I|}\delta(A_t).
 \tag{COV.2}
\]
This proves positivity for every nonempty $U$. If
$B\leq A$ and $U$ properly contains $F(B)$, every $A_t$
contains $B$ and at least one properly contains it.
Strict strongness gives
$\delta(U)>|I|\delta(B)=\delta(F(B))$.
Thus $F(B)\leq F(A)$. Inducedness and functoriality follow
directly from the definition.

The persistent family uses the nine-vertex object $W_*$
with order
\[
 a_0<a_1<a_2<b_0<b_1<b_2<c_0<c_1<c_2
\]
and five edges
\[
 a_0b_0c_0,\quad a_0b_1c_1,\quad a_0b_2c_2,
 \quad a_1b_0c_1,\quad a_2b_1c_2.
 \tag{COV.3}
\]
It has $\delta(W_*)=3$. It is linear, and the minimum
numbers of vertices used by $1,2,3,4,5$ of its edges are,
respectively, $3,5,6,8,9$. Hence every subset with at least
two vertices has predimension at least three. All nine
singletons are strong, and no other proper nonempty subset
is strong. The exact certificate
\texttt{code/C4\_voltage\_cover\_certificate.py}
checks all $512$ subsets and these edge-union bounds.

Let $a_0,a_1,a_2:V\to W_*$ denote the three singleton
arrows with the indicated images. Their marginal-difference
map is injective over every coefficient ring. Indeed, for
a finite collection of $r$ distinct strong copies of $W_*$
in a target, intersections of copies are strong in each
copy and so have at most one vertex. Their five-edge sets
are therefore pairwise disjoint. If $v$ is the number of
vertices appearing in the three selected positions, the
whole union has at most $v+6r$ vertices and at least $5r$
edges. Positivity gives
\[
 0<\delta(\text{union})\leq 2v-3r.
 \tag{COV.4}
\]
If every selected vertex belonged to at least two selected
triples, then $2v\leq3r$, a contradiction. Thus some selected
vertex belongs to exactly one copy. Its coefficient in at
least one of the two marginal differences is $1$ or $-1$.
Applying this argument to the nonzero support of a kernel
element forces a nonzero coefficient to vanish. This proves
injectivity over arbitrary nonzero rings, and excludes an
affine equalizer.

Fix $i\in I$. There is an increasing induced section
$s:W_*\to F(W_*)$ given by
\[
 s(a_j)=(a_j,i),\qquad s(b_j)=(b_j,\sigma(i)),
 \qquad s(c_j)=(c_j,\tau(i)).
 \tag{COV.5}
\]
It preserves all five edges and introduces none. It is
strong: for a subset $U$ properly containing its image,
one has $A_1=W_*$ and some nonempty $A_t$ with $t\geq2$.
Equation (COV.2) and positivity imply
$\delta(U)>\delta(W_*)$.
The singleton section $t:V\to F(V)$ choosing the index
$i$ is strong for the same reason. These sections obey
\[
 F(a_j)t=s a_j\qquad(j=0,1,2).
 \tag{COV.6}
\]
An affine equalizer of the three $F(a_j)$ could therefore
be composed on the right with $t$ and restricted through
$s$, producing an affine equalizer of the original three
arrows with unchanged augmentation. This is impossible
over every nonzero ring.

Composing two covers replaces the ordered fiber by the
lexicographically ordered product and the two permutations
by their componentwise products. Thus every finite
composition has the same form and the same obstruction.
In a filtered colimit of categories, any finite affine
witness, including its finitely many equality checks,
occurs at some finite stage. Applying this to the image
of (COV.3)'s family contradicts the obstruction just proved.

For completeness, filtered colimits preserve totality:
restrict a set action on the colimit category to all stages.
A finite simplicial sphere in a connected component of its
category of elements, together with finitely many paths
connecting its vertices, lifts to a common stage. That
stage's action component is contractible, so the sphere
has a nullhomotopy there and hence in the colimit.
The same finite-diagram argument handles components and
$\pi_1$. Since each stage $\Cdelta$ is total, the claimed
filtered colimit is total.
\end{proof}

\paragraph{Scope.}
The conclusion covers arbitrary finite fiber sizes and
arbitrary fixed permutations in (COV.1), not only cyclic
prime fibers. It does not cover arbitrary endofunctors,
templates that replace edges by larger internal gadgets,
or a modification in which the transition is not such a
finite cover. The total all-prime implication remains open.
